% =============================================================================
%   TYPHOON — Rapport de stage
%   Service de données climatiques par bâtiment pour l'assurance française
%
%   Stage d'été — SummerCamp 2026, Talan.  Auteur : Hazem Mrad.
%
%   Document court (20–30 pages) : le problème métier, la réponse technique,
%   l'architecture en diagrammes UML, les résultats mesurés, et les points de
%   douleur qui restent — chacun suivi de son remède.
%
%   Compilation :  tectonic -X compile main.tex
%   Les 13 chapitres de la version longue sont conservés dans archive-2026-10/
% =============================================================================
\documentclass[11pt,a4paper,oneside]{report}

% ── Portabilité du moteur ────────────────────────────────────────────────────
% pdfLaTeX a besoin de inputenc + fontenc ; XeTeX (tectonic, XeLaTeX, Overleaf
% en mode XeLaTeX) lit l'UTF-8 nativement et attend fontspec. Sans ce garde-fou,
% une compilation pdfLaTeX échoue sur les accents.
\ifdefined\XeTeXversion
  \usepackage{fontspec}
\else
  \usepackage[utf8]{inputenc}
  \usepackage[T1]{fontenc}
  \usepackage{lmodern}
\fi

\usepackage[french]{babel}
\usepackage{geometry}
\geometry{margin=2.1cm,top=2.4cm,bottom=2.3cm}

\usepackage{graphicx}
% Les captures d'écran vivent dans docs/screenshots/ (celles du README) :
% une seule copie dans le dépôt, référencée par son nom de fichier.
\graphicspath{{../docs/screenshots/}}
\usepackage{xcolor}
\usepackage{tikz}
\usetikzlibrary{positioning,arrows.meta,fit,backgrounds,calc,shapes.geometric,shapes.multipart,decorations.pathreplacing,chains}
\usepackage{booktabs}
\usepackage{array}
\usepackage{tabularx}
\usepackage{tcolorbox}
\tcbuselibrary{skins,breakable}
\usepackage{enumitem}
\usepackage{titlesec}
\usepackage{fancyhdr}
\usepackage{caption}
\usepackage{listings}
\usepackage{hyperref}

% ── Palette ──────────────────────────────────────────────────────────────────
\definecolor{typhoonDeep}{HTML}{0E2C3F}
\definecolor{typhoonBlue}{HTML}{1B4965}
\definecolor{typhoonAccent}{HTML}{2E7DAF}
\definecolor{typhoonAmber}{HTML}{B8791A}
\definecolor{typhoonRed}{HTML}{9E2B2B}
\definecolor{typhoonGreen}{HTML}{256B4E}
\definecolor{typhoonGrey}{HTML}{5A6472}
\definecolor{typhoonSand}{HTML}{F4F1EA}

% =============================================================================
%  MÉTADONNÉES DU RAPPORT — seul bloc à modifier pour la page de garde
% =============================================================================
\newcommand{\reportAuthor}{Hazem Mrad}
\newcommand{\reportSchool}{ESPRIT}
\newcommand{\reportSchoolLong}{École Supérieure Privée d'Ingénierie et de Technologie}
\newcommand{\reportDegree}{Diplôme d'ingénieur}
\newcommand{\reportCycle}{Cycle d'ingénieur}
\newcommand{\reportCompany}{Talan}
\newcommand{\reportCompanyLong}{Groupe Talan --- Conseil et ingénierie technologique}
\newcommand{\reportProgramme}{SummerCamp 2026}
\newcommand{\reportPeriod}{14 août 2026 -- 15 septembre 2026}
\newcommand{\reportAcademicYear}{2026--2027}
\newcommand{\reportTutorAcademic}{\dotfill}
\newcommand{\reportTutorCompany}{\dotfill}

% ── Logo de l'école ----------------------------------------------------------
% Déposer le fichier sous rapport/logos/esprit.png (.jpg ou .pdf aussi).
% Sans fichier, un cadre de remplacement explicite est affiché : le rapport
% compile quand même.
\newcommand{\reportSchoolLogo}[1][2.4cm]{%
  \IfFileExists{logos/esprit.png}{\includegraphics[height=#1]{logos/esprit.png}}{%
  \IfFileExists{logos/esprit.jpg}{\includegraphics[height=#1]{logos/esprit.jpg}}{%
  \IfFileExists{logos/esprit-logo.jpg}{\includegraphics[height=#1]{logos/esprit-logo.jpg}}{%
  \IfFileExists{logos/esprit-logo.png}{\includegraphics[height=#1]{logos/esprit-logo.png}}{%
  \IfFileExists{logos/esprit.pdf}{\includegraphics[height=#1]{logos/esprit.pdf}}{%
    \begin{tikzpicture}
      \node[draw=typhoonGrey,dashed,rounded corners=3pt,inner sep=9pt,
            text width=2.4cm,align=center,font=\sffamily\scriptsize,
            text=typhoonGrey,minimum height=#1]
        {Logo de l'école\\[2pt]à déposer dans\\[2pt]\texttt{logos/}};
    \end{tikzpicture}}}}}}%
}

% ── Logo Talan ───────────────────────────────────────────────────────────────
% Le tracé vectoriel `logos/talan-logo.tex` (blanc, extrait du SVG officiel)
% reste préféré : posé sur le bandeau sombre du titre, il est net à toute
% échelle et ne dépend d'aucun fichier binaire. L'image `logos/talan-logo.png`
% déposée dans le dossier est le repli. Aucun des deux ⇒ aucune erreur : c'est
% l'espace qui reste vide, jamais le document qui échoue.
\newcommand{\reportCompanyLogo}{%
  \IfFileExists{logos/talan-logo.tex}{%
    \input{logos/talan-logo}%
  }{%
    \IfFileExists{logos/talan-logo.png}{%
      \includegraphics[height=1.15cm]{logos/talan-logo.png}%
    }{}%
  }%
}

% ── Titres ───────────────────────────────────────────────────────────────────
\titleformat{\chapter}[display]
  {\normalfont\sffamily\bfseries\color{typhoonBlue}}
  {\filright\normalsize\color{typhoonGrey}\MakeUppercase{\chaptertitlename}~\thechapter}
  {6pt}
  {\Huge}
  [\vspace{2pt}{\color{typhoonAccent}\titlerule[1.1pt]}]
\titleformat{\section}{\normalfont\sffamily\large\bfseries\color{typhoonBlue}}{\thesection}{0.7em}{}
\titleformat{\subsection}{\normalfont\sffamily\normalsize\bfseries\color{typhoonGrey}}{\thesubsection}{0.6em}{}
\titlespacing*{\chapter}{0pt}{-18pt}{22pt}
\titlespacing*{\section}{0pt}{14pt}{6pt}
\titlespacing*{\subsection}{0pt}{10pt}{4pt}

% ── En-têtes ─────────────────────────────────────────────────────────────────
\pagestyle{fancy}
\fancyhf{}
\fancyhead[L]{\footnotesize\sffamily\color{typhoonGrey}\nouppercase{\leftmark}}
\fancyhead[R]{\footnotesize\sffamily\color{typhoonGrey}Typhoon}
\fancyfoot[C]{\footnotesize\sffamily\color{typhoonGrey}\thepage}
\setlength{\headheight}{22pt}
\renewcommand{\headrulewidth}{0.4pt}
\renewcommand{\headrule}{\hbox to\headwidth{\color{typhoonAccent}\leaders\hrule height \headrulewidth\hfill}}

% ── Listes compactes ─────────────────────────────────────────────────────────
\setlist[itemize]{topsep=2pt,itemsep=1.5pt,parsep=0pt,leftmargin=1.2em}
\setlist[enumerate]{topsep=2pt,itemsep=1.5pt,parsep=0pt,leftmargin=1.5em}

% ── Boîtes ───────────────────────────────────────────────────────────────────
\newtcolorbox{invariant}[1][]{colback=typhoonBlue!5,colframe=typhoonBlue,
  boxrule=0.5pt,left=6pt,right=6pt,top=4pt,bottom=5pt,breakable,
  fonttitle=\bfseries\sffamily\small,title={Règle non négociable},#1}
\newtcolorbox{douleur}[1][]{colback=typhoonRed!4,colframe=typhoonRed,
  boxrule=0.5pt,left=6pt,right=6pt,top=4pt,bottom=5pt,breakable,
  fonttitle=\bfseries\sffamily\small,title={Douleur},#1}
\newtcolorbox{remede}[1][]{colback=typhoonGreen!5,colframe=typhoonGreen,
  boxrule=0.5pt,left=6pt,right=6pt,top=4pt,bottom=5pt,breakable,
  fonttitle=\bfseries\sffamily\small,title={Remède},#1}
\newtcolorbox{encadre}[1][]{colback=typhoonSand,colframe=typhoonGrey!55,
  boxrule=0.4pt,left=6pt,right=6pt,top=4pt,bottom=5pt,breakable,
  fonttitle=\bfseries\sffamily\small,#1}

% ── Styles TikZ pour les diagrammes UML ──────────────────────────────────────
\tikzset{
  uml/.style={draw=typhoonBlue!75,fill=typhoonBlue!5,rounded corners=2pt,
              align=center,font=\small,inner sep=5pt},
  umlcomp/.style={uml,minimum width=3.0cm,minimum height=0.95cm,fill=typhoonBlue!8},
  umlactor/.style={uml,fill=typhoonAmber!15,rounded corners=0pt,minimum width=2.3cm},
  umllayer/.style={draw=typhoonAccent!45,dashed,rounded corners=3pt,inner sep=8pt,
                   fill=typhoonAccent!3},
  umlnote/.style={draw=black!35,fill=black!3,rounded corners=1pt,
                  align=left,font=\footnotesize,inner sep=4pt},
  umlstate/.style={uml,fill=typhoonGreen!8,rounded corners=8pt,minimum width=2.4cm},
  flow/.style={-{Latex[length=1.8mm]},draw=typhoonBlue!85,thick},
  flowd/.style={-{Latex[length=1.8mm]},draw=typhoonBlue!85,thick,dashed},
  flowback/.style={-{Latex[length=1.8mm]},draw=typhoonBlue!85,thick,dashed},
  flowwarn/.style={-{Latex[length=1.8mm]},draw=typhoonRed!80,thick,dashed},
  flowok/.style={-{Latex[length=1.8mm]},draw=typhoonGreen!80,thick},
}

\captionsetup{font=small,labelfont={bf,sf,color=typhoonBlue},skip=5pt}
\lstset{basicstyle=\ttfamily\footnotesize,breaklines=true,frame=single,
  rulecolor=\color{typhoonGrey!45},backgroundcolor=\color{typhoonSand},
  keywordstyle=\color{typhoonBlue},commentstyle=\color{typhoonGreen},
  stringstyle=\color{typhoonRed},showstringspaces=false,tabsize=2,
  xleftmargin=6pt,xrightmargin=6pt,aboveskip=6pt,belowskip=6pt}

\hypersetup{colorlinks=true,linkcolor=typhoonAccent,urlcolor=typhoonAccent,
  citecolor=typhoonAccent,
  pdftitle={Rapport de stage — Typhoon, service de données climatiques par bâtiment},
  pdfauthor={Hazem Mrad},
  pdfsubject={Rapport de stage — SummerCamp 2026, Talan}}

\setlength{\parindent}{0pt}
\setlength{\parskip}{4pt}
% Les figures du document sont hautes (diagrammes UML, captures d'écran) :
% avec les seuils par défaut, une figure de 12 cm ne tient ni en bas ni en haut
% d'une page de texte et part sur une page de flottants — trois pages blanches
% à moitié vides pour deux figures. Ces seuils autorisent une page pleine ; le
% document y gagne deux pages.
\renewcommand{\topfraction}{0.92}
\renewcommand{\bottomfraction}{0.92}
\renewcommand{\textfraction}{0.06}
\renewcommand{\floatpagefraction}{0.85}
% Tolérance de justification : évite les débordements de quelques points dans
% les légendes et les cellules étroites.
\setlength{\emergencystretch}{2em}

% ── Plafond de largeur des figures ───────────────────────────────────────────
% Un diagramme plus large que la justification serait rogné à la marge droite
% (LaTeX signale un « Overfull \hbox » et coupe). \figbox plafonne la figure à
% \textwidth sans jamais l'agrandir : un petit diagramme garde sa taille.
\newcommand{\figbox}[1]{%
  \resizebox{\ifdim\width>\textwidth\textwidth\else\width\fi}{!}{#1}%
}

% ── Captures d'écran ─────────────────────────────────────────────────────────
% \ecran{<fichier>}{<légende>} : une capture en demi-largeur, cadre fin, légende
% sous l'image. Deux par ligne via \hfill — les captures sont larges, une par
% ligne coûterait une page à elles seules.
\newlength{\ecranwidth}
\setlength{\ecranwidth}{0.485\textwidth}
\newcommand{\ecran}[2]{%
  \begin{minipage}[t]{\ecranwidth}
    \centering
    \setlength{\fboxsep}{0pt}\setlength{\fboxrule}{0.4pt}%
    % 1 pt de retrait : le filet du cadre s'ajoute sinon à la largeur et
    % déborde de la minipage (Overfull \hbox de 0,8 pt sur chaque capture).
    \fbox{\includegraphics[width=\dimexpr\linewidth-1pt\relax]{#1}}\\[3pt]
    {\scriptsize\sffamily\color{typhoonGrey}#2}
  \end{minipage}%
}
% Latin Modern n'a pas de petites capitales grasses sans-serif : la police est
% silencieusement remplacée. On déclare la substitution pour supprimer
% l'avertissement « Font shape TU/lmss/bx/sc undefined ».
\DeclareFontShape{TU}{lmss}{bx}{sc}{<->ssub*lmss/bx/n}{}

% Raccourcis
\newcommand{\code}[1]{{\ttfamily\small #1}}
\newcommand{\nom}[1]{\textsf{\small #1}}

% =============================================================================
\begin{document}

% -----------------------------------------------------------------------------
%  Page de garde — usages académiques français
% -----------------------------------------------------------------------------
\begin{titlepage}
\thispagestyle{empty}

% Bandeau supérieur : identité de l'entreprise d'accueil
\begin{tikzpicture}[remember picture,overlay]
  \fill[typhoonDeep] (current page.north west) rectangle
      ([yshift=-3.4cm]current page.north east);
  \fill[typhoonBlue] ([yshift=-3.4cm]current page.north west) rectangle
      ([yshift=-3.46cm]current page.north east);
\end{tikzpicture}

\noindent\begin{minipage}[c]{0.44\textwidth}
  \color{white}\sffamily
  {\bfseries\large \reportCompany}\\[2pt]
  {\footnotesize \reportProgramme}\\[2pt]
  {\scriptsize\color{white!75}\reportCompanyLong}
\end{minipage}\hfill
\begin{minipage}[c]{0.30\textwidth}
  \raggedleft\reportCompanyLogo
\end{minipage}

\vspace{1.3cm}

\begin{center}
  \reportSchoolLogo[2.2cm]\\[8pt]
  {\sffamily\bfseries\large \reportSchool}\\[2pt]
  {\sffamily\small\color{typhoonGrey} \reportSchoolLong}
\end{center}

\vspace{1.0cm}

\begin{center}
  {\sffamily\small\color{typhoonGrey}\MakeUppercase{Rapport de stage de fin de cycle}}\\[6pt]
  {\color{typhoonBlue}\rule{4cm}{0.8pt}}\\[12pt]

  {\sffamily\bfseries\LARGE\color{typhoonDeep}
   Service de données climatiques\\[4pt] par bâtiment}\\[10pt]

  {\sffamily\large\color{typhoonDeep}
   Concevoir et industrialiser une jointure aléa $\times$ vulnérabilité\\[2pt]
   à l'échelle du bâtiment pour les assureurs français}\\[8pt]

  {\sffamily\small\color{typhoonGrey}
   De la dispersion des données publiques à un contrat de données traçable,\\
   testé et déployé sur une infrastructure serverless.}
\end{center}

\vfill

\begin{center}
\renewcommand{\arraystretch}{1.4}
\begin{tabular}{@{}p{4.6cm}@{\hspace{4pt}}p{7.6cm}@{}}
  \textcolor{typhoonGrey}{\sffamily\small Présenté par} & \textbf{\reportAuthor} \\
  \textcolor{typhoonGrey}{\sffamily\small Intitulé du diplôme} & \reportDegree \\
  \textcolor{typhoonGrey}{\sffamily\small Cycle} & \reportCycle \\
  \textcolor{typhoonGrey}{\sffamily\small Structure d'accueil} & \reportCompanyLong \\
  \textcolor{typhoonGrey}{\sffamily\small Période de stage} & \reportPeriod \\
  \textcolor{typhoonGrey}{\sffamily\small Tuteur pédagogique} & \reportTutorAcademic \\
  \textcolor{typhoonGrey}{\sffamily\small Tuteur entreprise} & \reportTutorCompany \\
  \textcolor{typhoonGrey}{\sffamily\small Année universitaire} & \textbf{\reportAcademicYear} \\
\end{tabular}
\end{center}

\vfill

\begin{center}
  {\color{typhoonBlue}\rule{\textwidth}{0.6pt}}\\[4pt]
  {\sffamily\footnotesize\color{typhoonGrey}
   Rapport de stage --- \reportAcademicYear \hspace{2em}·\hspace{2em}
   \reportAuthor \hspace{2em}·\hspace{2em} \reportCompany}
\end{center}

\end{titlepage}

% ── Résumé ───────────────────────────────────────────────────────────────────
\chapter*{Résumé}
\addcontentsline{toc}{chapter}{Résumé}
\markboth{Résumé}{Résumé}

Évaluer le risque climatique d'un bien immobilier suppose aujourd'hui de croiser
des sources publiques dispersées, hétérogènes et sans clé de jointure commune.
L'assureur comme le promoteur paient cette dispersion en heures d'analyse, en
incertitude et en risques mal évalués.

Typhoon prend une adresse et rend un \textbf{enregistrement canonique} : les
aléas du référentiel Géorisques confrontés au point (et non à la commune), la
vulnérabilité du bâtiment issue de la BDNB (139 champs repris verbatim), la
classe de crue réglementaire TRI, l'hydrographie, la météo et la photo de
terrain. Chaque fait déclare sa provenance et son niveau de résolution ; une
absence de donnée est un fait, jamais une valeur par défaut.

Ce document présente le problème métier (\S1 à \S2), l'architecture en
diagrammes UML (\S3 à \S7), les résultats mesurés (\S8), puis les points de
douleur qui subsistent et leur remède (\S9).

\begin{encadre}[title={Ce que le document ne prétend pas}]
Typhoon ne produit ni score composite, ni bande de risque, ni recommandation
générée : le modèle actuariel reste celui de l'assureur. Cette limite est
délibérée et structurante — elle est imposée par le schéma du contrat, qui
\textbf{refuse} ces champs à la sérialisation.
\end{encadre}

\clearpage

% ── Sommaire ─────────────────────────────────────────────────────────────────
\renewcommand{\contentsname}{Sommaire}
\tableofcontents

\clearpage

% ── Chapitres ────────────────────────────────────────────────────────────────
\chapter{Introduction}
\label{ch:introduction}

\section{Le risque climatique est devenu une ligne de compte}

Un assureur dommages français ne peut plus tarifer un bien immobilier sans
regarder ce qui lui arrive : submersion, ruissellement, mouvements de terrain,
retrait-gonflement des argiles, séisme, feux de forêt. Un promoteur ne peut plus
livrer un programme sans répondre à la même question, posée par son assureur,
son banquier ou son acquéreur.

Les données existent. elles sont publiques, souvent de grande qualité, et
produites par des acteurs identifiés : Géorisques (BRGM, MTE), la BDNB (CSTB),
le géocodage de la Géoplateforme (IGN), l'hydrographie de la BD TOPO, Open-Meteo,
Panoramax. Le problème n'est donc pas l'absence de données : c'est leur
\textbf{dispersion}, leur \textbf{hétérogénéité} et l'absence de clé de jointure
commune.

\begin{encadre}[title={Le décalage central}]
Les référentiels publics décrivent le risque à l'échelle de la \textbf{commune}
ou du \textbf{secteur}. Une décision d'assurance ou d'investissement se prend à
l'échelle du \textbf{bâtiment}. Tout l'écart entre les deux est du travail manuel
— et, quand il est franchi sans méthode, une source d'erreur silencieuse.
\end{encadre}

\section{Objet du document}

Ce rapport décrit \textbf{Typhoon}, un service dont la fonction est de combler
cet écart : à partir d'une adresse, produire un enregistrement canonique qui
joint les aléas du référentiel Géorisques au bâtiment lui-même, en déclarant
pour chaque fait son origine et sa finesse.

Le document est organisé en trois mouvements, dans l'ordre du problème vers la
solution :

\begin{itemize}
  \item \textbf{le problème métier} (\S\ref{ch:douleurs}) — les douleurs que
    subissent concrètement l'assureur et le promoteur, formulées en coûts et en
    risques, et non en fonctionnalités ;
  \item \textbf{la réponse} (\S\ref{ch:reponse} à \S\ref{ch:robustesse}) — le
    contrat de données, l'architecture en diagrammes UML, les parcours réels et
    les garanties de robustesse ;
  \item \textbf{le bilan} (\S\ref{ch:resultats} et \S\ref{ch:douleurs-restantes})
    — ce qui a été mesuré, puis les points de douleur qui \emph{subsistent},
    chacun suivi de son remède.
\end{itemize}

\section{Les acteurs}

Quatre profils utilisent le service, avec des besoins différents qui expliquent
la forme du contrat de données.

\begin{itemize}
  \item \textbf{L'assureur} — tarifie un portefeuille, doit justifier sa
    décision devant son régulateur et ne veut pas d'une boîte noire. Il lui faut
    la provenance de chaque fait et le droit de ne pas conclure quand la source
    n'a pas répondu.
  \item \textbf{Le promoteur} — engage un programme immobilier et doit répondre
    vite à une question de risque, en amont de l'acquisition.
  \item \textbf{Le gestionnaire de portefeuille} — suit un ensemble de biens
    dans le temps et compare ce qui est comparable.
  \item \textbf{L'analyste} — ouvre la fiche bâtiment, lit les champs BDNB
    verbatim et cherche ce qui motive un écart de tarification.
\end{itemize}

\section{Périmètre et non-périmètre}

\begin{tabularx}{\textwidth}{@{}lX@{}}
\toprule
\textbf{Dans le périmètre} & \textbf{Hors périmètre (par décision)} \\
\midrule
Diagnostic par bâtiment, France métropolitaine & Calcul de prime, tarification actuarielle \\
Aléas Géorisques + classe TRI au point & Score composite, bande de risque, notation \\
Vulnérabilité BDNB verbatim (139 champs) & Recommandation générée présentée comme un fait \\
Provenance, résolution, attribution sur chaque fait & Prédiction de sinistre \\
Estimation de dommages paramétrique (scénarios TRI) & Données Copernicus dans le contrat sérialisé \\
\bottomrule
\end{tabularx}

\begin{invariant}
\textbf{Aucun score, aucune bande, aucune recommandation générée} n'entre dans
le contrat sérialisé. Ce n'est pas une réserve commerciale, c'est une contrainte
technique : le schéma canonique du backend
(\code{diagnostic\_record.py}) \emph{refuse} ces champs à la validation. Une
donnée qui n'a pas de source ne peut pas être ajoutée au contrat sans qu'un
amendement explicite soit voté.
\end{invariant}

\chapter{Le problème : ce que subissent assureurs et promoteurs}
\label{ch:douleurs}

Ce chapitre décrit les douleurs telles qu'elles se manifestent dans le travail
réel, formulées en \textbf{coûts}, en \textbf{délais} et en \textbf{risques
juridiques}. Chacune est ensuite reprise en \S\ref{ch:reponse} sous la forme
« douleur \(\rightarrow\) réponse technique ».

\section{Sept douleurs}

\subsection{D1 — L'adresse n'est pas l'unité des données publiques}

Les aléas sont publiés à la commune ou par emprises cartographiques ; le
bâtiment, lui, est décrit par la BDNB. Répondre « au bâtiment » avec une donnée
communale produit une réponse qui paraît précise et ne l'est pas.

\begin{douleur}[title={D1 — Granularité inadaptée}]
\textbf{Effet} : sur-estimation sur les biens situés hors des emprises, ou
sous-estimation sur ceux situés à l'intérieur d'une emprise communale large.
\textbf{Coût} : reprise manuelle pour les cas litigieux, et perte de confiance
dans l'outil dès le premier écart détecté par un expert.
\end{douleur}

\subsection{D2 — Aucune clé de jointure commune}

Le géocodeur rend une adresse et un point ; la BDNB rend un
\code{batiment\_groupe\_id} et une géométrie ; Géorisques rend des polygones
d'aléas ; l'hydrographie rend des tronçons ; la photo rend une prise de vue. Il
n'existe pas d'identifiant partagé : la jointure se fait par
\textbf{géométrie} (point dans polygone) et par codes administratifs, avec toutes
les approximations que cela implique.

\subsection{D3 — Résolution hétérogène, non déclarée}

Certaines données sont communales (l'aléa retrait-gonflement des argiles), la
plupart sont vectorielles et ponctuelles, quelques-unes sont surfaciques. Mélanger
ces niveaux sans le dire revient à présenter une estimation communale comme une
mesure sur le bien.

\begin{invariant}[title={Chaque fait déclare sa finesse}]
Tout objet d'aléa du contrat porte \code{source}, \code{resolution} et
l'attribution LO 2.0 du producteur. L'aléa argile est \emph{toujours}
\code{commune-level-estimate} tant qu'aucune source vectorielle n'est câblée :
la changer exige un amendement de la constitution du projet.
\end{invariant}

\subsection{D4 — « Absence de donnée » confondue avec « absence de risque »}

C'est la douleur la plus coûteuse, parce qu'elle est \emph{invisible}. Une source
qui répond 500, un quota atteint, un service momentanément indisponible : si
l'outil rend alors « aucun aléa », l'utilisateur conclut à un bien sans risque.
Le pire des deux mondes — une décision prise sur du vide, avec l'apparence d'une
réponse.

\subsection{D5 — Traçabilité insuffisante pour justifier une décision}

Un assureur doit pouvoir expliquer une tarification. Une valeur sans producteur,
sans date de millésime et sans méthode n'est pas défendable devant un régulateur
ni devant un client qui conteste.

\subsection{D6 — Le coût et le délai de l'analyse manuelle}

Croiser six sources par bien, à la main, représente typiquement plusieurs heures
par dossier. Cela rend l'analyse exhaustive d'un portefeuille économiquement
impossible : on échantillonne, donc on subit le risque de queue.

\subsection{D7 — La dette juridique de l'IA générative}

Introduire un modèle de langage dans un document réglementaire sans garde-fou
revient à accepter qu'un chiffre inventé soit publié. La douleur n'est pas
technique, elle est juridique : la responsabilité du contenu reste à celui qui le
signe.

\section{Les douleurs, mises en regard}

\begin{table}[h]
\centering
\small
\begin{tabularx}{\textwidth}{@{}p{4.2cm}Xp{3.4cm}@{}}
\toprule
\textbf{Douleur} & \textbf{Manifestation observable} & \textbf{Ce qu'elle coûte} \\
\midrule
D1 Granularité & Écarts entre l'outil et l'expert & Crédibilité \\
D2 Jointure & Scripts ad hoc, cas non couverts & Maintenance \\
D3 Résolution & Un chiffre communal présenté comme un fait du bien & Risque juridique \\
D4 Absence \(\neq\) 0 & Décision prise sur une source muette & Sinistre non anticipé \\
D5 Traçabilité & Valeur sans producteur ni millésime & Non-conformité \\
D6 Volume & Analyse limitée à un échantillon & Risque de queue \\
D7 IA non bornée & Chiffre plausible mais faux & Responsabilité \\
\bottomrule
\end{tabularx}
\caption{Les sept douleurs, leur manifestation et ce qu'elles coûtent réellement.}
\end{table}

\section{Le test de réalité : un portefeuille de 1 200 biens}

Pour rendre le problème tangible, prenons un cas représentatif.

Un assureur doit réévaluer un portefeuille de 1 200 biens répartis sur le
territoire. Le travail d'analyse disponible représente l'équivalent de deux
postes à temps plein. À plusieurs heures par bien, seule une fraction du
portefeuille peut être instruite en profondeur ; le reste est traité par
approximation communale.

Trois conséquences s'enchaînent :

\begin{enumerate}
  \item \textbf{Sur les biens non instruits}, la tarification retenue est celle
    de la commune — par exemple une parcelle en bord de cours d'eau traitée
    comme une parcelle quelconque du fait d'une emprise mal comprise.
  \item \textbf{Sur les biens contestés}, l'analyse doit être refaite à la main
    pour produire une justification acceptable, sans capitalisation du travail
    déjà fait.
  \item \textbf{Sur l'ensemble}, aucune réponse n'est reproductible : deux
    analystes, deux méthodes, deux résultats.
\end{enumerate}

\begin{encadre}[title={Ce que la réponse doit donc garantir}]
Une réponse \emph{par bâtiment}, \emph{traçable}, \emph{reproductible}, et
capable de dire « je ne sais pas » sans le confondre avec « il n'y a rien ».
C'est exactement le cahier des charges de \S\ref{ch:reponse}.
\end{encadre}

\chapter{La réponse}
\label{ch:reponse}

\section{La promesse en une phrase}

\begin{center}
\large\sffamily\color{typhoonBlue}
Une adresse entre ; un \textbf{enregistrement canonique par bâtiment} sort, où
chaque fait porte sa source, sa finesse et son attribution.
\end{center}

\begin{figure}[h]
\centering
\figbox{
\begin{tikzpicture}[node distance=1.5cm and 1.2cm]
  \node[uml,minimum width=2.6cm,fill=typhoonAmber!15] (in) {Une adresse\\\footnotesize (texte libre)};
  \node[umlcomp,right=of in,minimum width=3.4cm,fill=typhoonBlue!10] (core)
    {\textbf{Chaîne de collecte}\\\footnotesize 6 sources officielles};
  \node[uml,right=of core,minimum width=3.4cm,fill=typhoonGreen!10] (out)
    {\textbf{Contrat canonique}\\\footnotesize aléas + BDNB + provenance};
  \draw[flow] (in) -- (core);
  \draw[flow] (core) -- (out);
  \node[umlnote,below=0.9cm of core,text width=7.6cm]
    {Jamais de valeur par défaut : une source muette produit une \emph{absence}
     déclarée, pas un zéro. Le schéma refuse les champs interdits
     (score, bande, recommandation générée).};
  \draw[flowd] (core) -- (out |- core.south);
\end{tikzpicture}
}
\caption{La chaîne, de l'adresse au contrat. La provenance est portée par le
contrat lui-même, pas rajoutée après coup.}
\label{fig:promesse}
\end{figure}

\section{Chaque douleur reçoit une réponse nommée}

\begin{table}[h]
\centering
\small
\begin{tabularx}{\textwidth}{@{}p{3.6cm}p{5.2cm}X@{}}
\toprule
\textbf{Douleur} & \textbf{Réponse} & \textbf{Garantie vérifiée} \\
\midrule
D1 Granularité & Jointure \textbf{au point} : le polygone d'aléa est intersecté
  avec les coordonnées exactes de l'adresse & Le résultat d'aléa est
  \code{point-in-polygon}, jamais un report communal \\
D2 Jointure & Pipeline unique côté serveur : géocodage, puis jointure par
  géométrie et par identifiants administratifs & Une seule implémentation ;
  le lot est \(N \times\) le pipeline unitaire \\
D3 Résolution & Chaque objet d'aléa déclare \code{resolution} &
  L'argile reste \code{commune-level-estimate} \\
D4 Absence \(\neq\) 0 & Une panne de source devient une \emph{absence déclarée},
  distincte d'une valeur nulle & Cinq états d'absence distincts et nommés
  (\S\ref{sec:etats-tri}) \\
D5 Traçabilité & \code{source} + millésime + attribution LO 2.0 sur chaque fait &
  Champs obligatoires du schéma \\
D6 Volume & Lot = \(N \times\) pipeline unitaire, avec budgets et quotas
  protégés côté client & Budget mensuel BDNB ; limite de débit sur les routes
  de collecte \\
D7 IA non bornée & La prose ne peut pas introduire un fait : elle est
  \textbf{validée} avant publication, champ par champ & Un nombre non ancré
  fait rejeter la phrase et revenir au gabarit calculé \\
\bottomrule
\end{tabularx}
\caption{Correspondance douleur \(\rightarrow\) réponse \(\rightarrow\) garantie.}
\end{table}

\section{Les invariants qui structurent tout le reste}

Quatre règles expliquent la majorité des décisions techniques décrites en
\S\ref{ch:architecture} et \S\ref{ch:donnees}.

\begin{invariant}[title={I1 — Pas de jugement dans le contrat}]
Aucun score, aucune bande D03, aucune recommandation générée ne peut être
sérialisé. Le modèle actuariel reste celui de l'assureur. Le schéma l'impose
techniquement, pas seulement par convention.
\end{invariant}

\begin{invariant}[title={I2 — L'absence est un fait, pas une panne}]
« Hors TRI » est un fait sur le territoire ; « service indisponible » est un
échec de source ; « dans un TRI sans classe au point » est un troisième état.
Ces états ont des formulations distinctes et ne se confondent jamais avec zéro.
\end{invariant}

\begin{invariant}[title={I3 — Un seul ancrage pour les chiffres}]
Le validateur qui garde la prose du rapport est \emph{le même} que celui qui
garde l'assistant documentaire : il est importé, jamais dupliqué. Un chiffre
absent du rapport fait refuser la réponse, il n'est jamais affiché.
\end{invariant}

\begin{invariant}[title={I4 — Une seule implémentation du calcul}]
Le traitement par lot n'est pas un second moteur : c'est \(N\) appels au
pipeline unitaire. Toute optimisation de volume qui créerait un second chemin de
calcul est refusée par construction.
\end{invariant}

\section{Le parcours utilisateur, en quatre étapes}

L'application suit le déroulé du travail d'analyse, et chaque étape a une seule
responsabilité.

\begin{enumerate}
  \item \textbf{Diagnostic} — saisie de l'adresse, géocodage, collecte des aléas
    et de la fiche bâtiment. C'est la seule étape qui interroge les sources.
  \item \textbf{Scénario} — choix réglementaire : la \emph{classe de crue TRI}
    fixe le niveau d'eau. La pluie prévue n'en fixe que la forme dans le temps ;
    elle ne pilote jamais l'intensité.
  \item \textbf{Jumeau} — représentation 3D des volumes du bâti réel, avec les
    figures des aléas réellement constatés au point. Le bien est mis en évidence
    par son identité (son identifiant BDNB), jamais par un plus-proche-voisin.
  \item \textbf{Rapport} — document déterministe, exportable en PDF, complété par
    une prose bornée et par un assistant documentaire qui refuse de répondre hors
    des faits du rapport.
\end{enumerate}

\begin{encadre}[title={Un mode explicitement séparé}]
Le jumeau comporte un \textbf{bac à sable} d'expérience (hypothèses, sans valeur
réglementaire). Il est entré par un contrôle explicite, il n'écrit rien — aucun
champ, aucun rapport, aucun PDF — et sa dépendance physique est chargée
dynamiquement : le chemin de données factuel ne paie rien pour lui. Un test dédié
vérifie qu'aucun module factuel ne l'importe.
\end{encadre}

\section{Ce que la réponse ne fait pas}

\begin{itemize}
  \item Elle ne calcule pas de prime et ne propose pas de tarif.
  \item Elle ne classe pas les biens entre eux et ne produit aucune notation.
  \item Elle n'invente pas de magnitude : quand aucune grandeur réelle ne pilote
    un aléa, seule sa \emph{présence} est représentée.
  \item Elle ne garantit pas la reproductibilité de la \emph{prose} : seuls le
    rapport calculé, ses nombres et ses actions le sont.
\end{itemize}

\chapter{Architecture}
\label{ch:architecture}

\section{Vue d'ensemble}

L'architecture suit une contrainte simple : \textbf{une seule implémentation du
calcul}, et une frontière explicite autour du contrat de données. Le schéma de
la figure~\ref{fig:composants} en montre les couches.

\begin{figure}[h]
\centering
\figbox{
\begin{tikzpicture}[node distance=0.45cm]

  % ── Présentation ──
  \node[umlcomp] (spa) {SPA React + Vite\\\footnotesize carte Mapbox 3D, panneaux, rapport};
  \node[umlcomp,right=0.5cm of spa] (twin) {Jumeau 3D\\\footnotesize environnement WebGL dédié};
  \node[umllayer,fit=(spa)(twin),
        label={[font=\footnotesize\sffamily\color{typhoonGrey}]above:Présentation (navigateur)}] (lPres) {};

  % ── API ──
  \node[umlcomp,minimum width=11.6cm,below=0.75cm of lPres] (api)
    {\textbf{API FastAPI} — \footnotesize routes \code{/diagnostic}, \code{/api/flood-alea},
     \code{/api/hydro}, \code{/api/meteo}, \code{/api/photo}, \code{/api/report}, \code{/health}};
  \node[umllayer,fit=(api),
        label={[font=\footnotesize\sffamily\color{typhoonGrey}]above:Exposition}] (lApi) {};

  % ── Services ──
  \node[umlcomp,minimum width=2.7cm,below=0.8cm of lApi.south west,anchor=north west] (canon)
    {pipeline canonique\\\footnotesize \code{build\_diagnostic\_record}};
  \node[umlcomp,minimum width=2.7cm,right=0.28cm of canon] (batch)
    {lot, budgets, limites\\\footnotesize \(N\times\) unitaire, quotas};
  \node[umlcomp,minimum width=2.7cm,right=0.28cm of batch] (agent)
    {rapport\\\footnotesize gabarit + prose validée};
  \node[umlcomp,minimum width=2.7cm,right=0.28cm of agent] (assist)
    {assistant ancré\\\footnotesize refuse hors périmètre};
  \node[umllayer,fit=(canon)(batch)(agent)(assist),
        label={[font=\footnotesize\sffamily\color{typhoonGrey}]above:Métier (services)}] (lServ) {};

  % ── Connecteurs ──
  \node[uml,minimum width=2.35cm,below=0.8cm of lServ.south west,anchor=north west] (c1) {BDNB};
  \node[uml,minimum width=2.35cm,right=0.24cm of c1] (c2) {Géorisques\\\footnotesize REST + WFS};
  \node[uml,minimum width=2.35cm,right=0.24cm of c2] (c3) {flood\_alea\\\footnotesize classes TRI};
  \node[uml,minimum width=2.35cm,right=0.24cm of c3] (c4) {hydro\\\footnotesize Sandre / BD TOPO};
  \node[uml,minimum width=2.35cm,below=0.22cm of c1] (c5) {météo\\\footnotesize Open-Meteo};
  \node[uml,minimum width=2.35cm,right=0.24cm of c5] (c6) {géocodage\\\footnotesize BAN / Géoplateforme};
  \node[uml,minimum width=2.35cm,right=0.24cm of c6] (c7) {photo\\\footnotesize Panoramax};
  \node[umllayer,fit=(c1)(c2)(c3)(c4)(c5)(c6)(c7),
        label={[font=\footnotesize\sffamily\color{typhoonGrey}]above:Connecteurs (une source = un module)}] (lConn) {};

  % ── Externes ──
  \node[uml,minimum width=4.2cm,below=0.7cm of lConn.south] (llm)
    {Fournisseur de prose\\\footnotesize Gemini (défaut) ou Mistral};
  \node[umlnote,right=0.5cm of llm,text width=7.4cm]
    {Réglable par configuration. Une valeur inconnue ne choisit
     \emph{aucun} fournisseur : le rapport sort en gabarit plutôt que vers un
     tiers non choisi.};

  % ── Flux ──
  \draw[flow] (lPres.south) -- node[right,font=\footnotesize]{HTTPS / JSON} (lApi.north);
  \draw[flow] (lApi.south) -- (lServ.north);
  \draw[flow] (lServ.south) -- (lConn.north);
  \draw[flow] (lConn.south) -- (llm.north);
  \draw[flowd] (agent.south) to[out=-90,in=90] (llm.north west);

\end{tikzpicture}
}
\caption{Diagramme de composants. Les connecteurs sont isolés des services : un
service ne parle jamais directement à une source.}
\label{fig:composants}
\end{figure}

\section{Les couches et leur responsabilité}

\begin{table}[h]
\centering
\small
\begin{tabularx}{\textwidth}{@{}p{3.5cm}p{5.6cm}X@{}}
\toprule
\textbf{Couche} & \textbf{Contenu} & \textbf{Responsabilité unique} \\
\midrule
Présentation & React, Vite, TypeScript, Mapbox & Afficher des faits ; ne calcule
  aucun risque \\
Exposition & FastAPI, schémas Pydantic & Valider l'entrée, refuser les champs
  interdits, limiter le débit \\
Métier & \code{canonical}, \code{batch}, budgets, rapport, assistant & Calculer,
  assembler, valider la prose \\
Connecteurs & Une source par module & Parler à une source, normaliser sa
  réponse, ne rien inventer \\
Externe & Prose LLM & Reformuler uniquement, sous validation \\
\bottomrule
\end{tabularx}
\caption{Le découpage en couches et la responsabilité unique de chacune.}
\end{table}

\section{Le contrat comme frontière}

La pièce centrale n'est pas une base de données : c'est un \textbf{schéma}. Le
contrat canonique est validé par Pydantic, et son intérêt est autant ce qu'il
accepte que ce qu'il \emph{refuse}.

\begin{lstlisting}[language=Python,caption={Extrait du schéma : les champs interdits sont refusés, pas ignorés.}]
class DiagnosticRecord(BaseModel):
    # ... adresse, geometrie, aleas, batiment, provenance ...
    model_config = ConfigDict(extra="forbid")   # un champ inconnu = erreur

    # Explicitement ABSENTS du contrat :
    #   niveau            -> bande D03 : display-only, hors contrat
    #   score             -> jugement fusionne, interdit (constitution 2)
    #   recommandations   -> narration IA, interdite
    #   copernicus        -> hors perimetre v1
\end{lstlisting}

\begin{invariant}
Un champ interdit ne peut pas être « ajouté discrètement » : il faut modifier le
schéma, donc passer par un amendement explicite. C'est ce qui rend la règle
tenable dans le temps — elle n'est pas une consigne d'équipe, elle est une
propriété du code.
\end{invariant}

\section{Déploiement}

\begin{figure}[h]
\centering
\figbox{
\begin{tikzpicture}[node distance=0.5cm and 1.1cm]

  % ── Poste utilisateur ──
  \node[uml,minimum width=4.0cm] (poste) {\textbf{Poste utilisateur}\\\footnotesize navigateur};
  \node[uml,minimum width=3.6cm,below=0.18cm of poste] (spa) {\footnotesize SPA (bundle statique)};

  % ── Hébergement ──
  \node[umlcomp,minimum width=4.4cm,right=of poste] (vercel)
    {\textbf{Option A — Vercel}\\\footnotesize fonction Python + statique};
  \node[umlcomp,minimum width=4.4cm,below=0.55cm of vercel] (render)
    {\textbf{Option B — Render}\\\footnotesize uvicorn \code{:8000}};
  \node[umlnote,below=0.4cm of render,text width=4.4cm]
    {Cache de rapports en mémoire de processus : \textbf{non persistant} en
     serverless (recalcul à chaque appel).};

  % ── Externe ──
  \node[uml,minimum width=4.2cm,right=of vercel,fill=typhoonAmber!10] (ext)
    {\textbf{Sources officielles}\\\footnotesize Géorisques, BDNB, IGN,\\Open-Meteo, Panoramax};
  \node[uml,minimum width=4.2cm,below=0.55cm of ext,fill=typhoonGreen!10] (llm)
    {\textbf{Fournisseur de prose}\\\footnotesize API Generative Language};

  \draw[flow] (poste) -- node[above,font=\footnotesize]{HTTPS} (vercel);
  \draw[flow] (render) -- (vercel);
  \draw[flow] (vercel) -- node[above,font=\footnotesize]{HTTPS} (ext);
  \draw[flow] (render.east) to[out=0,in=180] node[below,font=\footnotesize]{HTTPS} (ext.west);
  \draw[flow] (render) -- (llm);
\end{tikzpicture}
}
\caption{Diagramme de déploiement. Les deux options d'hébergement exposent la
même API ; seule la persistance du cache diffère.}
\label{fig:deploiement}
\end{figure}

Deux mécanismes transverses protègent l'ensemble :

\begin{itemize}
  \item \textbf{La limitation de débit} s'applique aux routes qui déclenchent une
    vraie collecte, et non aux lectures légères : les chemins sont énumérés
    exactement, ce qui évite d'attraper par préfixe les routes qui n'ont pas à
    être freinées.
  \item \textbf{Les quotas amont sont protégés côté client} : la BDNB est
    appelée sous un plafond de 120 requêtes par minute et le géocodage sous
    50 par seconde, avec une réattente unique sur un 429 explicite.
\end{itemize}

\section{Les choix d'architecture, et pourquoi}

\begin{table}[h]
\centering
\small
\begin{tabularx}{\textwidth}{@{}p{4.6cm}X@{}}
\toprule
\textbf{Choix} & \textbf{Raison} \\
\midrule
FastAPI + Pydantic & Le contrat est le livrable ; le valider à l'entrée et à la
  sortie est plus utile qu'un schéma documentaire \\
Un module par source & Une source qui tombe ne doit pas contaminer les autres, et
  un parseur par format (REST, GML 3.2) reste lisible \\
Lot \(= N \times\) unitaire & Deux moteurs divergeraient : un chiffre pour
  l'unité, un autre pour le portefeuille \\
Client HTTP dédié au WFS & Le \emph{keep-alive} du WFS Géorisques est lié à son
  backend ; le mélanger avec les appels REST produit des 404 en cascade \\
Prose validée \emph{après} génération & Aucun prompt ne garantit qu'un modèle ne
  produira pas un chiffre plausible ; la garantie est dans le validateur \\
\bottomrule
\end{tabularx}
\caption{Cinq décisions structurantes et leur justification.}
\end{table}

\chapter{Le modèle de données}
\label{ch:donnees}

\section{Diagramme de classes}

Deux grappes de classes cohabitent : le \textbf{contrat de collecte}, qui
représente ce que les sources ont répondu, et le \textbf{contrat du rapport}, qui
représente ce que le document affirme. Aucune classe du second n'invente un
champ qui n'existe pas dans le premier.

\begin{figure}[h]
\centering
\figbox{
\begin{tikzpicture}[
  cls/.style={draw=typhoonBlue!75,fill=typhoonBlue!4,rounded corners=1.5pt,
              inner sep=2.5pt,font=\scriptsize},
  rel/.style={draw=typhoonBlue!80,thick},
  font=\small]

  % ── Contrat de collecte ──
  \node[cls] (dr) {\begin{tabular}{@{}l@{}}
    \textbf{DiagnosticRecord}\\\hline
    adresse : string\\
    geometrie : Point (lat, lon)\\
    aleas : AléaObjet [0..13]\\
    batiment : BdnbFiche [0..1]\\
    hydro : HydroModel [0..1]\\
    provenance : BlocSource\\\hline
    validation : refuse les champs interdits
  \end{tabular}};

  \node[cls,below=0.75cm of dr] (alea) {\begin{tabular}{@{}l@{}}
    \textbf{AléaObjet}\\\hline
    code : string\\
    present : boolean\\
    \textbf{source} : string\\
    \textbf{resolution} : enum\\
    \textbf{attribution} : string\\
    intensity? : seulement si\\
    \ \ une grandeur réelle la pilote
  \end{tabular}};

  \node[cls,right=0.6cm of alea] (bdnb) {\begin{tabular}{@{}l@{}}
    \textbf{BdnbFiche}\\\hline
    batiment\_groupe\_id : string\\
    champs : 139 verbatim\\
    empreinte : géométrie\\
    hauteur, altitude, usage
  \end{tabular}};

  % ── Contrat du rapport ──
  \node[cls,right=1.4cm of dr] (rr) {\begin{tabular}{@{}l@{}}
    \textbf{ReportRequest}\\\hline
    sector : string\\
    scenario : ScenarioModel\\
    damage : DamageEstimate\\
    timestamp : string
  \end{tabular}};

  \node[cls,below=0.7cm of rr] (sc) {\begin{tabular}{@{}l@{}}
    \textbf{ScenarioModel}\\\hline
    key : extreme | moyen |\\
    \ \ frequent | faible \ ($=$ TRI)\\
    risk : string\\
    depthPeakM : float
  \end{tabular}};

  \node[cls,right=0.45cm of sc] (dmg) {\begin{tabular}{@{}l@{}}
    \textbf{DamageEstimate}\\\hline
    damageEUR : Range\\
    floodedConduitM : Range\\
    damagedBuildings : Range\\
    \dots\ (10 indicateurs)
  \end{tabular}};

  \node[cls,below=0.7cm of dmg] (asst) {\begin{tabular}{@{}l@{}}
    \textbf{AssistantResponse}\\\hline
    answer : string\\
    answered\_from\_report : bool\\
    refusal : out\_of\_scope |\\
    \ \ refused\_number | unavailable\\
    model : string
  \end{tabular}};

  % ── Relations ──
  \draw[rel] (dr) -- node[right,font=\scriptsize]{1} node[left,font=\scriptsize]{0..13} (alea);
  \draw[rel] (dr.east) -- node[above,font=\scriptsize]{0..1} (rr.west);
  \draw[rel] (alea) -- node[above,font=\scriptsize]{0..1} (bdnb);
  \draw[rel] (rr) -- node[left,font=\scriptsize]{1} (sc);
  \draw[rel] (rr) -- node[left,font=\scriptsize]{1} (dmg);
  \draw[rel,flowd] (rr.east) to[out=0,in=90] node[right,font=\scriptsize]{assistant ancré sur \emph{ce} rapport} (asst.north);

  \node[umlnote,below=0.5cm of bdnb,text width=4.3cm]
    {\textbf{Jamais dans le contrat :} \code{niveau}/\code{score},
     \code{recommandations} générées, \code{copernicus}. Le schéma les refuse.};
\end{tikzpicture}
}
\caption{Diagramme de classes du contrat. Les trois champs en gras d'\code{AléaObjet}
sont obligatoires : un aléa sans provenance n'est pas représentable.}
\label{fig:classes}
\end{figure}

\section{La provenance n'est pas une option}

Chaque objet d'aléa porte trois champs qui rendent le fait vérifiable par
l'utilisateur lui-même.

\begin{itemize}
  \item \code{source} — le producteur du fait (par exemple « Géorisques (BRGM /
    MTE) »), jamais le nom de notre module interne.
  \item \code{resolution} — la finesse réelle : \code{point-in-polygon},
    \code{building-level} ou \code{commune-level-estimate}.
  \item \code{attribution} — la mention de licence ouverte (LO 2.0), à afficher.
\end{itemize}

\begin{lstlisting}[caption={Un objet d'aléa tel qu'il sort du pipeline (extrait, valeurs réelles).}]
{
  "code": "inondation",
  "present": true,
  "source": "Georisques (BRGM / MTE)",
  "resolution": "point-in-polygon",
  "attribution": "Source Georisques -- donnees a jour au <millesime>",
  "evidence": { "in_tri": true, "classe": "extreme",
                "water_depth_m": 1.1, "origin": "tri-class" }
}
\end{lstlisting}

\begin{invariant}[title={L'argile reste une estimation communale}]
L'aléa retrait-gonflement des argiles est \textbf{toujours} publié en
\code{commune-level-estimate} tant qu'aucune source vectorielle n'est câblée. La
règle est écrite dans la constitution du projet : la changer exige un
amendement, pas un commit.
\end{invariant}

\section{La fiche bâtiment, verbatim}

La BDNB fournit 139 champs décrivant le bâtiment. Ils sont repris
\textbf{verbatim}, sans interprétation : matériaux de mur, de toiture, année de
construction, hauteur, altitude au sol, usage, empreinte. Le service n'ajoute ni
normalisation ni catégorisation.

Ce choix a un coût assumé : le contrat est volumineux, et certaines valeurs sont
vides. Il a un bénéfice direct : un expert peut vérifier chaque valeur à la
source, et aucune perte d'information n'est introduite par notre traitement. Là
où la BDNB ne répond pas, le contrat porte \code{null} — jamais une valeur
plausible inventée.

\section{Le niveau d'eau vient d'un seul endroit}

Le jumeau numérique affiche une hauteur d'eau. Cette valeur provient d'un unique
module d'entrée et n'est calculée nulle part ailleurs. En particulier, le modèle
d'impact ne doit jamais importer ce module : sans cette contrainte, une seconde
définition du niveau d'eau apparaîtrait, et deux affichages du même rapport
finiraient par diverger.

\begin{invariant}[title={Le scénario fixe le niveau, la pluie n'en fixe que la forme}]
La \emph{classe de crue TRI officielle} détermine le pic. La prévision de pluie ne
détermine que la courbe dans le temps. Hors TRI, la hauteur reste strictement
nulle. Une journée sèche produit une rampe explicitement étiquetée comme
hypothèse — elle ne fait jamais disparaître le scénario.
\end{invariant}

\section{Volumétrie}

\begin{table}[h]
\centering
\small
\begin{tabular}{@{}lrl@{}}
\toprule
\textbf{Objet} & \textbf{Volume} & \textbf{Remarque} \\
\midrule
Aléas du référentiel & 13 & toujours présents, y compris absents \\
Fiche bâtiment BDNB & 139 champs & verbatim, \code{null} si vide \\
Indicateurs d'impact & 10 & chacun en triple (valeur, basse, haute) \\
Sections du rapport & 5 & synthèse, catégories, mitigations, confiance, annexe \\
États d'absence TRI & 5 & distincts, nommés, non confondus (\S\ref{sec:etats-tri}) \\
\bottomrule
\end{tabular}
\caption{Volumétrie du contrat. Les 13 aléas sont toujours présents : c'est
l'absence qui est signifiante.}
\end{table}

\chapter{Les parcours}
\label{ch:parcours}

\newcommand{\acteur}[3]{%
  \draw[thick,typhoonBlue!80] (#1,#2) circle (2.2mm);
  \draw[thick,typhoonBlue!80] (#1,#2-0.22) -- (#1,#2-0.66);
  \draw[thick,typhoonBlue!80] (#1-0.30,#2-0.38) -- (#1+0.30,#2-0.38);
  \draw[thick,typhoonBlue!80] (#1,#2-0.66) -- (#1-0.26,#2-1.00);
  \draw[thick,typhoonBlue!80] (#1,#2-0.66) -- (#1+0.26,#2-1.00);
  \node[font=\footnotesize\sffamily,align=center,below] at (#1,#2-1.05) {#3};%
}

\section{Diagramme de cas d'usage}

Quatre acteurs, neuf cas d'usage. Le périmètre du système est explicite : tout ce
qui est en dehors relève de l'assureur (son modèle actuariel) et n'est pas
implémenté ici.

\begin{figure}[h]
\centering
\figbox{
\begin{tikzpicture}[uc/.style={ellipse,draw=typhoonAccent!80,fill=typhoonAccent!5,
    align=center,font=\scriptsize,minimum width=3.6cm,minimum height=0.82cm,inner sep=1pt}]

  % ── Frontière du système ──
  \draw[typhoonBlue!45,dashed,rounded corners=4pt] (2.9,0.1) rectangle (13.4,7.0);
  \node[font=\scriptsize\sffamily\color{typhoonGrey},anchor=south west] at (2.95,7.0)
    {Typhoon — service de données};

  % ── Acteurs ──
  \acteur{1.0}{6.1}{Assureur}
  \acteur{1.0}{3.5}{Gestionnaire\\de portefeuille}
  \acteur{1.0}{0.9}{Promoteur}
  \acteur{14.6}{3.5}{Analyste}

  % ── Cas d'usage ──
  \node[uc] (u1) at (6.1,6.2) {Diagnostiquer une adresse};
  \node[uc] (u2) at (6.1,5.0) {Consulter la fiche BDNB};
  \node[uc] (u3) at (6.1,3.8) {Choisir un scénario TRI};
  \node[uc] (u4) at (6.1,2.6) {Simuler dans le jumeau};
  \node[uc] (u5) at (6.1,1.4) {Générer le rapport};
  \node[uc] (v1) at (10.6,6.2) {Analyser un portefeuille};
  \node[uc] (v2) at (10.6,5.0) {Exporter le PDF};
  \node[uc] (v3) at (10.6,3.8) {Interroger l'assistant};
  \node[uc] (v4) at (10.6,2.6) {Suivre un bien dans le temps};

  % ── Associations ──
  \draw[flow] (1.5,6.1) -- (u1.west);
  \draw[flow] (1.5,6.1) -- (u3.west);
  \draw[flow] (1.5,3.5) -- (u5.west);
  \draw[flow] (1.5,3.5) -- (v1.west);
  \draw[flow] (1.5,0.9) -- (u4.west);
  \draw[flow] (1.5,0.9) -- (u5.west);
  \draw[flow] (14.1,3.5) -- (u2.east);
  \draw[flow] (14.1,3.5) -- (v3.east);
  \draw[flow] (u4.east) -- (v4.west);
  \draw[flow] (u5.east) -- (v2.west);
\end{tikzpicture}
}
\caption{Diagramme de cas d'usage. Le classement des biens et le calcul de
prime sont \emph{absents} du périmètre : ce sont les cas d'usage que Typhoon
refuse volontairement de couvrir.}
\label{fig:usecase}
\end{figure}

\section{Séquence 1 — diagnostiquer une adresse}

C'est le parcours qui interroge les sources. Il doit tenir deux promesses
simultanées : ne pas attendre indéfiniment, et ne jamais transformer une panne en
absence de risque. Le traitement par lot réutilise exactement ce parcours :
\(N\) appels au pipeline unitaire, aucun second moteur.

\begin{figure}[h]
\centering
\figbox{
\begin{tikzpicture}[font=\scriptsize,
    life/.style={draw=typhoonGrey!55,dashed},
    head/.style={umlcomp,minimum width=2.2cm,minimum height=0.7cm,font=\scriptsize}]

  % ── Participants ──
  \node[head] (nav) at (1.3,0) {Navigateur};
  \node[head] (api) at (4.6,0) {API FastAPI};
  \node[head] (geo) at (7.8,0) {Géocodeur BAN};
  \node[head] (conn) at (10.7,0) {Connecteurs};
  \node[head] (pipe) at (13.6,0) {Pipeline};

  \foreach \p in {nav,api,geo,conn,pipe} {\draw[life] (\p.south) -- ++(0,-6.1);}

  % ── Messages ──
  \draw[flow] (1.3,-0.6) -- node[above]{POST /diagnostic/adresse} (4.6,-0.6);
  \draw[flow] (4.6,-1.3) -- node[above]{géocodage} (7.8,-1.3);
  \draw[flowback] (7.8,-2.0) -- node[above]{point + codes INSEE} (4.6,-2.0);
  \draw[flow] (4.6,-2.7) -- node[above]{collecte parallèle (6 sources)} (10.7,-2.7);
  \draw[flowback] (10.7,-3.4) -- node[above]{faits bruts, ou absence déclarée} (4.6,-3.4);
  \draw[flow] (4.6,-4.1) -- node[above]{\code{build\_diagnostic\_record}} (13.6,-4.1);
  \draw[flowback] (13.6,-4.8) -- node[above]{contrat validé, champs interdits refusés} (4.6,-4.8);
  \draw[flowback] (4.6,-5.5) -- node[above]{200 — JSON canonique + provenance} (1.3,-5.5);

  % ── Notes, en bandeau sous le diagramme ──
  \node[umlnote,anchor=north west,text width=8.0cm] at (0.2,-6.4)
    {\textbf{429} : une seule réattente sur un refus explicite, puis échec franc.
     Le géocodeur est la source la plus sollicitée.};
  \node[umlnote,anchor=north west,text width=5.9cm] at (8.9,-6.4)
    {\textbf{Quotas} protégés côté client : BDNB 120/min, Géoplateforme 50/s.};
\end{tikzpicture}
}
\caption{Séquence d'un diagnostic. La réponse finale ne peut pas contenir « aucun
aléa » à la place d'une source en panne : les deux sont produits par des
branches distinctes.}
\label{fig:seq-diagnostic}
\end{figure}

\section{Séquence 2 — la prose du rapport}

Le rapport est \emph{template-first} : les sections calculées partent
immédiatement, et la prose arrive ensuite, si elle passe la validation. Le
document n'attend jamais le modèle pour exister.

\begin{figure}[h]
\centering
\figbox{
\begin{tikzpicture}[font=\scriptsize,
    life/.style={draw=typhoonGrey!55,dashed},
    head/.style={umlcomp,minimum width=2.6cm,minimum height=0.7cm,font=\scriptsize}]

  \node[head] (cli) at (1.4,0) {Navigateur (SSE)};
  \node[head] (str) at (5.2,0) {\code{report\_agent}};
  \node[head] (val) at (9.0,0) {Validateur d'ancrage};
  \node[head] (llm) at (13.0,0) {Fournisseur de prose};

  \foreach \p in {cli,str,val,llm} {\draw[life] (\p.south) -- ++(0,-6.3);}

  \draw[flow] (1.4,-0.6) -- node[above]{POST /api/report/stream} (5.2,-0.6);
  \draw[flowback] (5.2,-1.4) -- node[above]{sections calculées} (1.4,-1.4);
  \draw[flow] (5.2,-2.2) -- node[above]{prose demandée (JSON strict)} (13.0,-2.2);
  \draw[flowback] (13.0,-3.0) -- node[above]{JSON — ou rien (clé absente, erreur)} (5.2,-3.0);
  \draw[flow] (5.2,-3.8) -- node[above]{chaque champ, un par un} (9.0,-3.8);
  \draw[flowwarn] (9.0,-4.5) -- node[above]{rejeté : nombre non ancré} (5.2,-4.5);
  \draw[flowok] (9.0,-5.2) -- node[above]{accepté} (5.2,-5.2);
  \draw[flowback] (5.2,-5.9) -- node[above]{« patch », puis « done »} (1.4,-5.9);

  \node[umlnote,anchor=north west,text width=7.4cm] at (0.2,-6.6)
    {Le repli est \textbf{par champ} : un champ rejeté ne fait pas tomber les
     autres. Le modèle annoncé est celui du fournisseur actif.};
  \node[umlnote,anchor=north west,text width=6.6cm] at (8.2,-6.6)
    {\code{fallback\_used} dit, sans ambiguïté, si le document contient de la
     prose ou seulement le gabarit calculé.};
\end{tikzpicture}
}
\caption{Séquence de la prose. Le modèle est un rédacteur sous contrainte, pas
une source de faits : tout ce qu'il écrit est vérifié avant d'être publié.}
\label{fig:seq-prose}
\end{figure}

\section{Les écrans du parcours}

Les deux séquences précédentes décrivent le système ; les captures qui suivent
montrent ce qu'il rend. Elles ont été prises sur l'application en service, en
rejouant le parcours complet depuis une recherche d'adresse — pas composées pour
l'illustration : les chiffres, les classes et les états de source visibles sur
ces images sont ceux que la source a rendus à cet instant. La provenance de
l'environnement de capture est donnée au chapitre~\ref{ch:resultats}.

\begin{figure}[htbp]
\centering
\ecran{06-etape1-diagnostic.png}{\textbf{Étape 1 — Diagnostic.} Les aléas du
  référentiel au point, chacun avec son niveau de résolution ; la commune
  retenue et ses chiffres INSEE ; la carte 3D du bâti.}
\hfill
\ecran{07-etape2-scenarios.png}{\textbf{Étape 2 — Scénarios.} Les classes de
  crue TRI officielles ; la hauteur d'eau au point en découle, elle n'est jamais
  choisie par un curseur.}\\[8pt]
\ecran{08-etape3-jumeau.png}{\textbf{Étape 3 — Jumeau.} Le bâti en volume dans
  son propre environnement WebGL : panneau des faits à gauche, catalogue des
  catastrophes à droite, horloge de l'expérience en bas.}
\hfill
\ecran{09-rapport.png}{\textbf{Étape 4 — Rapport.} Page autonome : scénario de
  référence, faits clés, provenance. La prose estampillée est un appoint, jamais
  la source des nombres.}
\caption{Le tunnel, de l'adresse au rapport. Chaque étape n'ajoute pas des
chiffres : elle change ce qu'on peut \emph{faire} de ceux de l'étape
précédente.}
\label{fig:ecrans-tunnel}
\end{figure}

\begin{figure}[htbp]
\centering
\ecran{03-tableau-de-bord.png}{\textbf{Tableau de bord.} Ce que l'assureur
  regarde en arrivant : volume diagnostiqué, biens à expertiser, dernières
  analyses, et la qualité des faits.}
\hfill
\ecran{04-portefeuille.png}{\textbf{Portefeuille.} Le livre des actifs, trié par
  résolution du fait : vérifié à la parcelle, estimé à la commune, ou
  manquant.}
\caption{Les deux écrans d'entrée de l'assureur. Tous deux lisent le cache
local des diagnostics : ils ne rejouent aucune requête vers les sources.}
\label{fig:ecrans-assureur}
\end{figure}

\section{Ce que ces deux séquences garantissent}

\begin{itemize}
  \item \textbf{Le calcul précède la rédaction.} Tous les nombres, niveaux et
    actions de mitigation sont produits par du code déterministe. Le modèle
    réécrit, il ne décide pas.
  \item \textbf{La validation est unique.} Le même validateur garde la prose du
    rapport et les réponses de l'assistant documentaire : il est importé, jamais
    recopié.
  \item \textbf{L'échec est un état, pas une exception.} Une clé absente, un
    quota atteint, une sortie non-JSON conduisent tous à un rapport complet et
    annoncé comme tel.
\end{itemize}

\chapter{Robustesse : la dégradation honnête}
\label{ch:robustesse}

Un service qui dépend de six sources publiques tombe souvent — non pas
complètement, mais \emph{partiellement}. La qualité de l'architecture se mesure
donc moins à son fonctionnement nominal qu'à sa façon de se dégrader.

\section{Diagramme d'activité : collecter sans mentir}

\begin{figure}[h]
\centering
\figbox{
\begin{tikzpicture}[font=\scriptsize,
    act/.style={draw=typhoonBlue!75,fill=typhoonBlue!5,rounded corners=2pt,
                align=center,minimum height=0.75cm,inner sep=3pt},
    bar/.style={fill=black!75,minimum height=0.09cm},
    dec/.style={diamond,draw=typhoonAmber!90,fill=typhoonAmber!8,aspect=2.2,
                align=center,inner sep=1pt}]

  \node[circle,fill=black!75,minimum size=3mm,inner sep=0pt] (start) at (4.5,0.4) {};
  \node[bar,minimum width=7.8cm] (fork) at (4.5,-0.35) {};

  \node[act,minimum width=1.85cm] (a1) at (0.9,-1.3) {BDNB};
  \node[act,minimum width=1.85cm] (a2) at (2.9,-1.3) {Géorisques};
  \node[act,minimum width=1.85cm] (a3) at (4.9,-1.3) {TRI au point};
  \node[act,minimum width=1.85cm] (a4) at (6.9,-1.3) {Hydro, météo};
  \node[act,minimum width=2.2cm] (a5) at (3.9,-2.25) {Photo de terrain};

  \node[bar,minimum width=7.8cm] (join) at (4.5,-2.95) {};

  \node[dec,minimum width=2.5cm,minimum height=1.0cm] (d) at (4.5,-3.85) {source\\muette ?};

  \node[act,minimum width=3.5cm,fill=typhoonRed!6,draw=typhoonRed!70] (abs) at (1.7,-5.15)
    {absence déclarée\\\scriptsize avec sa raison nommée};
  \node[act,minimum width=3.5cm,fill=typhoonGreen!7,draw=typhoonGreen!70] (fait) at (7.3,-5.15)
    {fait publié\\\scriptsize source + résolution + attribution};

  \node[act,minimum width=6.4cm] (asm) at (4.5,-6.35) {assembler et valider le contrat};
  \node[circle,draw=black!75,thick,minimum size=4.2mm,inner sep=0pt] (end) at (4.5,-7.15) {};
  \node[circle,fill=black!75,minimum size=2.4mm,inner sep=0pt] at (4.5,-7.15) {};

  \draw[flow] (start) -- (fork);
  \foreach \a in {a1,a2,a3,a4} {\draw[flow] (fork.south -| \a.north) -- (\a.north);}
  \draw[flow] (a2.south) -- (a5.north);
  \foreach \a in {a1,a2,a3,a4,a5} {\draw[flow] (\a.south) -- (\a.south |- join.north);}
  \draw[flow] (join) -- (d);
  \draw[flow] (d.west) -- node[above]{oui} (abs.north);
  \draw[flow] (d.east) -- node[above]{non} (fait.north);
  \draw[flow] (abs.south) -- (abs.south |- asm.north);
  \draw[flow] (fait.south) -- (fait.south |- asm.north);
  \draw[flow] (asm) -- (end);
\end{tikzpicture}
}
\caption{Diagramme d'activité de la collecte. La branche « source muette »
produit un \emph{fait} sur la source, pas une valeur sur le risque : c'est cette
distinction qui rend la dégradation exploitable.}
\label{fig:activite}
\end{figure}

\section{Diagramme d'états : les cinq vérités de l'absence}
\label{sec:etats-tri}

Le cas du périmètre TRI est le plus instructif, parce que « il n'y a pas de
classe de crue ici » recouvre en réalité cinq situations différentes. Les
confondre revient à présenter une panne comme une bonne nouvelle.

\begin{figure}[h]
\centering
\figbox{
\begin{tikzpicture}[font=\scriptsize,
    st/.style={draw=typhoonBlue!75,fill=typhoonBlue!5,rounded corners=9pt,
               align=center,minimum width=3.1cm,minimum height=0.8cm,inner sep=3pt},
    stok/.style={st,draw=typhoonGreen!75,fill=typhoonGreen!7},
    stwarn/.style={st,draw=typhoonRed!75,fill=typhoonRed!6}]

  \node[circle,fill=black!75,minimum size=3mm,inner sep=0pt] (init) at (0,1.5) {};
  \node[st] (q) at (0,0.6) {Interrogation du\\périmètre TRI};

  \node[stok] (hors) at (-2.6,-1.3) {\textbf{Hors TRI}\\fait territorial};
  \node[st] (sans) at (2.6,-1.3) {\textbf{Dans un TRI}\\sans classe au point};
  \node[stwarn] (nu) at (-2.6,-3.1) {\textbf{Appartenance}\\non vérifiée};
  \node[stwarn] (panne) at (2.6,-3.1) {\textbf{Service}\\indisponible};

  \draw[flow] (init) -- (q);
  \draw[flow] (q.south) -- ++(0,-0.25) -| (hors.north);
  \draw[flow] (q.south) -- ++(0,-0.25) -| (sans.north);
  \draw[flow] (q.south) -- ++(0,-0.5) -| (nu.north);
  \draw[flow] (q.south) -- ++(0,-0.5) -| (panne.north);

  \node[font=\tiny,anchor=east,align=right] at (-0.15,-0.62) {point hors emprise};
  \node[font=\tiny,anchor=west,align=left] at (0.15,-0.62) {polygone couvrant,\\ aucune hauteur publiée};
  \node[font=\tiny,anchor=east,align=right] at (-0.4,-1.95) {réponse inexploitable};
  \node[font=\tiny,anchor=west,align=left] at (0.4,-1.95) {erreur réseau, 5xx, quota};

  \node[umlnote,below=0.35cm of nu,text width=8.6cm,anchor=north]
    {Chaque état a sa \emph{formulation propre}, partagée entre le panneau de
     scénarios et le rapport (« ce n'est pas une absence de risque »). Aucun
     chemin ne produit « 0 m » à la place d'une source muette.};
\end{tikzpicture}
}
\caption{Diagramme d'états de l'absence TRI. « Hors TRI » est un fait ; « service
indisponible » est une panne : deux messages, deux traitements.}
\label{fig:etats-tri}
\end{figure}

\begin{invariant}[title={Absence \(\neq\) panne}]
Le connecteur porte cinq vérités distinctes : chargement, panne, dans un TRI sans
classe au point, hors TRI, appartenance non vérifiée. La formulation est unique et
partagée. Un panneau qui afficherait « pas de risque » sur une panne violerait la
règle la plus importante du projet.
\end{invariant}

\section{La discipline de test}

La robustesse n'est pas vérifiée par relecture mais par exécution. La suite
offline couvre 261 tests et ne dépend d'aucun accès réseau : toutes les sources
sont simulées par un transport HTTP factice, ce qui rend les scénarios de panne
\textbf{reproductibles} — un 500, un 429, un GML tronqué se rejouent à volonté.

\begin{itemize}
  \item \textbf{Les tests précèdent le code.} Un test qui échoue est commité
    avant l'implémentation qui le fait passer.
  \item \textbf{Les tests existants ne sont pas modifiés} pour faire passer une
    fonctionnalité : ce serait déplacer la barre au lieu de franchir l'obstacle.
  \item \textbf{Un garde-fou d'isolation} vérifie qu'aucun module factuel
    n'importe le bac à sable, et que le bac à sable n'importe ni le moteur de
    dommages, ni le rapport, ni la carte, ni le schéma du contrat.
  \item \textbf{La limitation de débit est testée} sous, à et au-delà du seuil,
    y compris le fait que les routes de lecture légère ne sont jamais freinées.
\end{itemize}

\section{Ce que la dégradation produit, en pratique}

Lors de la vérification en conditions réelles, deux sources ont effectivement
échoué pendant la même exécution : la BDNB a répondu un statut de limitation sur
sa première page, et le service TRI s'est déclaré partiellement indisponible.

Le résultat obtenu est exactement celui attendu par la figure~\ref{fig:activite} :

\begin{itemize}
  \item la fiche bâtiment a affiché des champs vides plutôt que des valeurs
    inventées ;
  \item le panneau TRI a écrit « point non cartographié » et rappelé que ce
    n'était pas une absence de risque ;
  \item le rapport a signalé douze aléas « non vérifiés » ;
  \item \textbf{aucune erreur serveur n'a été produite} : le diagnostic a abouti,
    partiellement renseigné et entièrement honnête.
\end{itemize}

\begin{encadre}[title={Le critère de réussite n'est pas le taux de succès}]
Un diagnostic dégradé mais exact vaut mieux qu'un diagnostic complet et faux. Le
service préfère rendre 80 \% de faits vérifiés avec leurs manques déclarés,
plutôt que 100 \% de champs remplis dont une partie serait devinée.
\end{encadre}

\chapter{Résultats mesurés}
\label{ch:resultats}

Tous les chiffres de ce chapitre proviennent d'exécutions réelles sur une
installation locale (API FastAPI, frontend Vite, sources publiques interrogées en
direct). Rien n'est extrapolé : ce qui n'a pas pu être mesuré est signalé comme
tel.

\begin{encadre}[title={Environnement des captures et des mesures}]
Les captures du chapitre~\ref{ch:parcours} ont été prises dans Chrome
(1440\,\texttimes\,900) sur le frontend servi par un dev server Vite local, la
session étant ouverte avec le compte de démonstration, et les adresses de
référence diagnostiquées au préalable pour alimenter le cache local qui nourrit
le tableau de bord et le portefeuille. L'API interrogée est celle de
l'\emph{instance déployée}, pour une raison d'environnement : le poste de capture
ne joint pas \code{www.georisques.gouv.fr}. Un diagnostic y reste bloqué jusqu'au
délai de chaque appel Géorisques (mesuré : 224\,s au lieu de 10\,s) et rend les
treize aléas en \code{erreur: source Géorisques indisponible} — la capture
n'aurait montré qu'un état de panne.

Deux conséquences sont visibles sur les images et ne sont pas des défauts de
l'application : l'instance déployée est en retard sur ce dépôt, donc
\code{/api/meteo/france} y répond 404 et les cinq cartes météo de l'étape 2
affichent « Référence météo indisponible » ; et l'adresse de référence rend les
aléas selon le millésime de l'instance, non celui du poste. Ce que le contrat
exige reste vérifiable sur ces captures : une source en panne produit un
\emph{état} nommé, jamais une valeur par défaut.
\end{encadre}

\section{Le diagnostic, en conditions réelles}

Un diagnostic sur « Quai de la Rapée, 75012 Paris » a produit les faits suivants.

\begin{table}[h]
\centering
\small
\begin{tabularx}{\textwidth}{@{}p{5.2cm}X@{}}
\toprule
\textbf{Élément observé} & \textbf{Résultat} \\
\midrule
Commune retenue & Paris 12\textsuperscript{e} arrondissement, INSEE 75112 (138 024 habitants, 8 431 hab./km²) \\
Aléas du référentiel & 13 objets dans le contrat, dont 4 constatés \emph{présents} au point (ICPE, inondation, sismicité, radon) \\
Classe de crue au point & dans le périmètre TRI ; classes \emph{moyen} et \emph{faible} cartographiées au point, hauteur d'eau réglementaire disponible \\
Hydrographie & réseau reconstruit côté serveur jusqu'à la mer (Seine) \\
Météo & prévision ponctuelle de l'adresse (et non moyenne départementale) \\
Photo de terrain & Panoramax, 24 candidates, la plus proche à 18,4\,m, orientée vers le bien \\
Erreur serveur & aucune \\
\bottomrule
\end{tabularx}
\caption{Résultat d'un diagnostic réel, exécuté pendant l'écriture de ce rapport.}
\end{table}

La durée d'un diagnostic complet se situe entre 9 et 16 secondes à froid sur
l'instance déployée, la première requête payant le réveil de la fonction ; la
même adresse rejouée depuis le cache local est instantanée. Quand une source
canonique est injoignable, la durée monte au délai cumulé des appels
(224 secondes mesurées) : c'est ce cas qui a justifié la \emph{borne} côté
client, plafonnée à cinq minutes et affichée à l'écran avec le temps écoulé.
À titre de comparaison, la reconstruction manuelle des mêmes faits demande
plusieurs heures par bien.

\begin{encadre}[title={Ce que cette mesure ne dit pas}]
Elle ne dit pas que l'analyse de risque prend 15 secondes : elle dit que
l'\emph{assemblage des faits} prend 15 secondes. L'interprétation, la
tarification et la décision restent humaines — c'est précisément la répartition
que le projet revendique.
\end{encadre}

\section{Le rapport}

\begin{table}[h]
\centering
\small
\begin{tabularx}{\textwidth}{@{}p{4.6cm}X@{}}
\toprule
\textbf{Scénario} & \textbf{Mesure} \\
\midrule
Sans fournisseur de prose configuré & 200 ; \code{fallback\_used = true} ; modèle annoncé \code{template-first} ; document complet de 2 752 caractères \\
Avec prose & 200 ; \code{fallback\_used = false} ; modèle annoncé \code{gemini-2.5-flash} ; 5 narratifs de catégorie et 4 narratifs de mitigation acceptés ; 2 850 caractères \\
Rejeu de la même requête & servi par le cache (\code{cache\_hit = true}), à l'identique \\
Flux SSE & séquence complète \code{header}, \code{summary}, \code{category} \(\times\) 5, \code{mitigations}, \code{confidence}, \code{appendix}, \code{done} \\
Flux SSE avec prose & 9 événements \code{patch} de prose réellement validée \\
\bottomrule
\end{tabularx}
\caption{Le rapport, mesuré avec et sans fournisseur de prose.}
\end{table}

\section{Le fournisseur de prose}

\begin{table}[h]
\centering
\small
\begin{tabular}{@{}ll@{}}
\toprule
\textbf{Opération} & \textbf{Durée mesurée} \\
\midrule
Appel structuré (JSON) & 1,1 seconde \\
Appel en flux (SSE) & 0,7 seconde \\
\bottomrule
\end{tabular}
\caption{Latence du fournisseur de prose sur des requêtes courtes.}
\end{table}

La vérification du câblage a été faite sans clé valide, par comparaison d'erreurs :
sans en-tête d'authentification, l'API répond
« \emph{Method doesn't allow unregistered callers} » ; avec un en-tête portant une
clé invalide, elle répond \code{API\_KEY\_INVALID}. Le \emph{changement} de
réponse prouve que l'en-tête est bien celui attendu — c'est un contrôle de
câblage, pas un test de bout en bout.

\section{La suite de tests}

\begin{table}[h]
\centering
\small
\begin{tabular}{@{}lll@{}}
\toprule
\textbf{Périmètre} & \textbf{Résultat} & \textbf{Durée} \\
\midrule
Backend, suite offline & 261 tests réussis, 0 échec (2 avertissements de ressource) & 3 min 54 \\
Frontend (Vitest) & 521 tests réussis, 0 échec, 63 fichiers & 28 s \\
Lint Python & aucun avertissement & — \\
Garde-fou d'isolation du bac à sable & vert & — \\
\bottomrule
\end{tabular}
\caption{État de la suite de tests au moment de la rédaction.}
\end{table}

Un point de méthode mérite d'être souligné : la suite reste \textbf{verte alors
qu'une clé d'API réelle est configurée} dans l'environnement. C'est une propriété
importante — un test qui appellerait le modèle en direct serait lent, non
déterministe et coûteux ; ici, toutes les sources sont simulées, et la présence
d'une clé ne change rien au résultat.

\section{Les quotas, protégés côté client}

\begin{table}[h]
\centering
\small
\begin{tabularx}{\textwidth}{@{}p{3.7cm}p{3.5cm}X@{}}
\toprule
\textbf{Source} & \textbf{Plafond appliqué} & \textbf{Comportement en cas de refus} \\
\midrule
BDNB & 120 requêtes / minute & budget mensuel de 10 000 appels ; retour partiel déclaré \\
Géocodage (Géoplateforme) & 50 requêtes / seconde & une seule réattente sur un 429 explicite, puis échec franc \\
Routes de collecte & 30 requêtes / 60 s & réponse 429 côté API, jamais un ralentissement silencieux \\
\bottomrule
\end{tabularx}
\caption{Les plafonds vérifiés chez le fournisseur et leur application.}
\end{table}

\chapter{Ce qui fait encore mal — et le remède}
\label{ch:douleurs-restantes}

\section{Deux lectures du même mot}

Le mot \emph{douleur} désigne ici deux choses qu'il faut distinguer :

\begin{itemize}
  \item la douleur \textbf{résolue} (\S\ref{ch:douleurs}) — ce que l'assureur et
    le promoteur subissaient avant, et que l'architecture traite par
    construction ;
  \item la douleur \textbf{restante} — ce qui, aujourd'hui, continue de faire mal
    à l'usage, et qui n'est pas caché dans ce rapport.
\end{itemize}

Cette seconde liste est la plus utile au lecteur : elle dit où le produit
n'est pas encore au niveau de sa promesse, et ce qu'il faudrait décider pour y
arriver.

\section{Les huit douleurs restantes}

\begin{table}[h]
\centering
\small
\begin{tabularx}{\textwidth}{@{}p{4.3cm}X@{}}
\toprule
\textbf{Douleur} & \textbf{Effet pour l'assureur ou le promoteur} \\
\midrule
P1 Prose refusée selon le format des nombres & Le résumé exécutif peut retomber
  en gabarit tandis que le reste est rédigé : document inégal devant un client \\
P2 Prose non reproductible & Deux demandes identiques peuvent produire deux
  formulations différentes \\
P3 Tier gratuit peu compatible avec des adresses réelles & Le contenu envoyé
  peut servir à l'amélioration des produits du fournisseur \\
P4 Communication publique en décalage & La page publique crédite encore un autre
  fournisseur et promet une souveraineté que le tier gratuit n'offre pas \\
P5 Aléa argile communal & Impossible de distinguer deux bâtiments d'une même
  commune — précisément sur l'aléa le plus sinistrogène \\
P6 Cache non persistant en serverless & Chaque rapport peut consommer un appel
  de modèle : un pic peut épuiser un quota et basculer en gabarit en pleine
  démonstration \\
P7 Dette de nommage & Le point d'entrée garde un nom de fournisseur qui n'est
  plus celui par défaut \\
P8 Outillage de test fragile & La suite ne démarre pas sans exclure un plugin
  tiers bloqué par la politique de sécurité de la machine \\
\bottomrule
\end{tabularx}
\caption{Les huit douleurs restantes, du point de vue de l'utilisateur.}
\end{table}

\section{Les trois douleurs qui coûtent le plus}

\subsection{P1 — La prose dépend du format des nombres}

C'est la douleur la plus visible, et elle a été \textbf{mesurée}. Le validateur
d'ancrage n'accepte que les nombres présents dans le rapport. Or le fournisseur
actuel écrit volontiers les nombres sous forme compacte ou à la française
(« 409 k€ », « 1,1 m »), là où l'ensemble autorisé ne contient que les formes
complètes (\code{408700}, \code{294264}, \code{523136}).

Résultat concret : les cinq narratifs de catégorie et les quatre narratifs de
mitigation sont acceptés, mais le \emph{résumé exécutif} est refusé et remplacé
par le gabarit calculé.

\begin{douleur}[title={Douleur observée}]
Le document mélange une prose de qualité et un premier paragraphe mécanique. Le
client ne connaît pas la cause et la lit comme un défaut de finition — alors
qu'il s'agit d'un garde-fou qui fonctionne.
\end{douleur}

\begin{remede}[title={Remède proposé}]
Élargir l'ensemble des nombres autorisés aux \textbf{faits du rapport lui-même} :
profondeur du scénario, bornes basses et hautes des dommages, et leurs formes
compactes (k€), plus les chiffres portés par le libellé du scénario. Le garde-fou
n'est pas affaibli : on autorise des faits du document, jamais des chiffres
extérieurs. Le changement exige un test sur une sortie réelle du modèle, car
cette autorité d'ancrage est partagée avec l'assistant documentaire.
\end{remede}

\subsection{P3 — Un tier gratuit qui n'est pas compatible avec des adresses}

Un assureur manipule des adresses de biens : la donnée est identifiante. Or, sur
le tier gratuit du fournisseur retenu par défaut, le contenu envoyé
\emph{peut servir à l'amélioration des produits} du fournisseur. Le même
connecteur, avec une clé facturée, ne présente pas cette réserve.

\begin{douleur}[title={Douleur}]
Le service fonctionne, mais pas dans les conditions qu'un client assurance
attendrait pour des données de portefeuille.
\end{douleur}

\begin{remede}[title={Remède proposé}]
Le fournisseur est déjà réglable par configuration : basculer sur une clé
facturée ne demande aucune modification de code. La décision à prendre est
commerciale (qui paie l'appel), pas technique. À défaut, un fournisseur
souverain peut être branché sans toucher au reste de l'architecture, puisque
l'interface du connecteur est neutralisée par un simple routage.
\end{remede}

\subsection{P5 — L'aléa argile reste communal}

C'est la douleur la plus coûteuse \emph{en sinistralité}. Le retrait-gonflement
des argiles est aujourd'hui publié en \code{commune-level-estimate} : deux
bâtiments d'une même commune reçoivent la même estimation, alors que leur
exposition réelle dépend du sol, de la végétation et du bâti.

\begin{douleur}[title={Douleur}]
Sur l'aléa qui pèse le plus dans les sinistres de sécheresse, le service ne
différencie pas les biens — ce qui est exactement ce qu'on lui demande par
ailleurs.
\end{douleur}

\begin{remede}[title={Remède proposé}]
Câbler une source vectorielle d'aléa argile et passer la résolution à
\code{building-level} ou \code{point-in-polygon}. La règle actuelle est inscrite
dans la constitution du projet : la modification doit donc passer par un
amendement explicite, puis par la spécification — jamais par une retouche de
code silencieuse.
\end{remede}

\section{Les cinq autres, en bref}

\begin{table}[h]
\centering
\small
\begin{tabularx}{\textwidth}{@{}p{1.1cm}X@{}}
\toprule
\textbf{} & \textbf{Douleur \(\rightarrow\) remède} \\
\midrule
P2 & Prose non reproductible \(\rightarrow\) le cache rend une même demande
  identique au rejeu ; pour le reste, la reproductibilité n'est pas promise, et
  aucun fait ne dépend de la prose. \\
P4 & Communication publique en décalage \(\rightarrow\) mettre la page publique,
  la FAQ et la mention légale en accord avec le fournisseur réellement utilisé,
  ou revenir à un fournisseur souverain. \\
P6 & Cache non persistant en serverless \(\rightarrow\) stockage externe, ou
  budget d'appels de prose journalier, sur le modèle du budget BDNB déjà en
  place. \\
P7 & Dette de nommage du point d'entrée \(\rightarrow\) renommage neutre dans un
  change set dédié, qui met à jour les tests qui épinglent ce chemin d'appel. \\
P8 & Outillage de test fragile \(\rightarrow\) exclusion documentée du plugin
  tiers en échec, ou image d'intégration continue maîtrisée, pour que
  « lancer les tests » redevienne une commande unique. \\
\bottomrule
\end{tabularx}
\caption{Les cinq douleurs restantes et leur remède.}
\end{table}

\section{Les décisions qui reviennent au commanditaire}

Trois questions ne sont pas techniques et ne peuvent pas être tranchées par
l'équipe :

\begin{enumerate}
  \item \textbf{Qui paie l'appel au modèle ?} Le tier gratuit est utilisable et
    suffisant en démonstration ; il n'est pas compatible avec des adresses de
    portefeuille réelles.
  \item \textbf{Où doit vivre la prose ?} En Europe (souveraineté) ou chez le
    fournisseur le moins cher (coût) — l'architecture supporte les deux, la
    décision est commerciale.
  \item \textbf{Jusqu'où automatiser l'interprétation ?} Le projet s'arrête
    volontairement avant le score et la recommandation. Aller plus loin exige un
    amendement de la constitution, pas un ajout de champ.
  \end{enumerate}

\chapter{Conclusion}
\label{ch:conclusion}

\section{Ce qui a été construit}

Typhoon transforme une adresse en un enregistrement canonique de risque par
bâtiment, en joignant les aléas du référentiel Géorisques au point exact et la
vulnérabilité réelle du bâti issue de la BDNB. Sept douleurs métier ont reçu
sept réponses vérifiables, et l'application les rend utilisables dans un parcours
en quatre étapes, du diagnostic au rapport exportable.

Le résultat tient moins à une prouesse technique qu'à une série de refus :

\begin{itemize}
  \item \textbf{refus du score} — le contrat ne peut pas porter de jugement, et
    le schéma l'empêche techniquement ;
  \item \textbf{refus du zéro par défaut} — une source muette produit une absence
    déclarée, jamais une valeur ;
  \item \textbf{refus du second moteur} — le lot reste \(N\) appels au pipeline
    unitaire ;
  \item \textbf{refus de la prose non vérifiée} — le modèle propose, le
    validateur dispose.
\end{itemize}

\section{Ce que le projet démontre}

Trois enseignements nous semblent transférables au-delà de ce service.

\begin{enumerate}
  \item \textbf{La provenance est une fonctionnalité, pas une métadonnée.} Dès
    qu'un fait est vérifiable par son destinataire, la discussion change de
    nature : elle porte sur la donnée, non sur la confiance dans l'outil.
  \item \textbf{Décider de ne pas savoir est une capacité.} Distinguer cinq
    états d'absence a demandé plus de travail que d'afficher zéro — et c'est ce
    qui rend le service utilisable dans un processus d'assurance.
  \item \textbf{Un modèle de langage devient utilisable dès qu'il est encadré.}
    Non par un prompt plus ferme, mais par un validateur qui compare ce qu'il
    écrit aux faits du document et refuse le reste.
\end{enumerate}

\section{Ce qui reste ouvert}

Le rapport a nommé huit douleurs restantes plutôt que de les passer sous
silence. Les trois plus lourdes sont l'ancrage des nombres selon leur
formatage (\S\ref{ch:douleurs-restantes}, P1), la compatibilité du tier gratuit
avec des adresses réelles (P3), et l'aléa argile qui demeure communal (P5) —
cette dernière exigeant un amendement de la constitution du projet avant toute
implémentation.

\begin{invariant}[title={La règle qui tient tout le reste}]
\textbf{Aucune donnée sans source ; aucune absence confondue avec un zéro ;
aucun chiffre non ancré publié.} Le reste de l'architecture n'est qu'une
manière de rendre ces trois phrases exécutables.
\end{invariant}


\appendix
\chapter{Annexe A — Surface d'API}
\label{ann:api}

Toutes les routes sont servies par l'API FastAPI. Les routes de \emph{collecte}
sont soumises à la limitation de débit ; les lectures légères ne le sont pas.

\begin{table}[h]
\centering
\small
\begin{tabularx}{\textwidth}{@{}p{6.3cm}p{1.6cm}X@{}}
\toprule
\textbf{Route} & \textbf{Méthode} & \textbf{Rôle} \\
\midrule
\code{/diagnostic/adresse} & POST & Adresse \(\rightarrow\) enregistrement canonique (pipeline unitaire) \\
\code{/diagnostic/batch} & POST & Soumission d'un lot (\(N \times\) unitaire) \\
\code{/diagnostic/batch/\{id\}} & GET & Suivi d'un lot (léger, non limité) \\
\code{/diagnostic/zone/building} & GET & Fiche BDNB par identifiant de bâtiment \\
\code{/diagnostic/zone/buildings} & GET & Bâtiments d'une emprise \\
\code{/diagnostic/adresse/rapport-pdf} & GET & Rapport exportable en PDF \\
\code{/api/flood-alea} & GET & Classes de crue TRI au point (point dans polygone) \\
\code{/api/hydro} & GET & Tracé hydrographique au point (Sandre / BD TOPO) \\
\code{/api/hydro/extend} & GET & Extension amont et aval du tracé \\
\code{/api/meteo} & GET & Pluie horaire et débit prévus (Open-Meteo) \\
\code{/api/meteo/france} & GET & Synthèse nationale (moyenne de points représentatifs) \\
\code{/api/photo} & GET & Photo de terrain et crédit (Panoramax) \\
\code{/api/geocode/search} & GET & Autocomplétion d'adresse (BAN / Géoplateforme) \\
\code{/api/report} & POST & Rapport complet (gabarit + prose validée) \\
\code{/api/report/stream} & POST & Même rapport en flux SSE (patch de prose) \\
\code{/api/report/assistant} & POST & Question sur \emph{ce} rapport, avec refus explicites \\
\code{/health} & GET & Sonde de vie \\
\code{/health/detailed} & GET & Sonde détaillée \\
\bottomrule
\end{tabularx}
\caption{Surface d'API. Les routes de collecte sont énumérées \emph{exactement}
dans le limiteur de débit, ce qui évite de freiner les lectures légères.}
\end{table}

\section{Codes et états}

\begin{itemize}
  \item \textbf{200} — la réponse est produite, éventuellement
    \emph{partiellement} renseignée (dégradation déclarée dans le corps).
  \item \textbf{422} — l'entrée ne satisfait pas le schéma (question vide ou trop
    longue, champ inconnu, \code{extra="forbid"}).
  \item \textbf{429} — plafond de débit atteint sur une route de collecte.
  \item \textbf{5xx} — réservé aux erreurs réellement inattendues : une source
    publique en panne n'en produit pas, elle devient une absence déclarée.
\end{itemize}

\section{Le cas de l'assistant documentaire}

La route \code{/api/report/assistant} mérite une mention : elle ne répond qu'à
partir des faits du rapport fourni dans la requête elle-même. Aucune recherche
web, aucune mémoire d'un rapport précédent, aucune persistance côté serveur — la
conversation n'existe que dans l'échange HTTP. Une question hors périmètre
obtient un \emph{refus} (un fait sur le document) ; une défaillance du modèle
obtient une \emph{indisponibilité} (une panne de source). Les deux messages sont
distincts, et jamais remplacés par une réponse inventée.

\chapter{Annexe B — Sources de données}
\label{ann:sources}

Chaque fait publié par le service provient d'une source identifiée. Ce tableau
est la version humaine de ce que le contrat porte champ par champ.

\begin{table}[h]
\centering
\small
\begin{tabularx}{\textwidth}{@{}p{2.9cm}p{3.3cm}p{2.6cm}X@{}}
\toprule
\textbf{Donnée} & \textbf{Producteur} & \textbf{Résolution} & \textbf{Usage dans le service} \\
\midrule
Aléas réglementaires & Géorisques (BRGM / MTE) & point dans polygone & présence et preuve de chaque aléa \\
Emprises de crue TRI & Géorisques (WFS) & point dans polygone & classe de crue et niveau d'eau du scénario \\
Fiche bâtiment & BDNB (CSTB) & bâtiment & 139 champs verbatim : matières, année, hauteur, empreinte \\
Adresse & BAN / Géoplateforme (IGN) & bâtiment (interpolation) & géocodage, point de jointure \\
Hydrographie & BD TOPO (IGN) / Sandre & tronçon & cours d'eau, distance, trajet de l'eau \\
Altitude & Altimétrie (IGN) & point & altitude au sol du bâtiment \\
Sol & Aléa argile (BRGM) & \textbf{commune} & sensibilité au retrait-gonflement \\
Météo prévue & Open-Meteo & point & pluie horaire, rafales, humidité \\
Vigilance crues & Vigicrues & tronçon & alerte en cours à proximité \\
Photo de terrain & Panoramax (IGN / DINUM) & prise de vue & contexte visuel et crédit \\
Prose du rapport & Fournisseur configurable & — & reformulation sous validation \\
\bottomrule
\end{tabularx}
\caption{Sources, producteurs et finesse réelle de chaque donnée.}
\end{table}

\section{Les règles de provenance}

\begin{itemize}
  \item Tout objet d'aléa porte \code{source}, \code{resolution} et
    l'attribution de licence ouverte (LO 2.0) du producteur.
  \item L'attribution est accompagnée d'un millésime : une donnée sans date de
    mise à jour n'est pas exploitable dans un contexte de conformité.
  \item Une source qui ne répond pas laisse une trace dans le contrat (absence
    déclarée) ; elle ne disparaît pas silencieusement.
  \item L'aléa argile est \textbf{toujours} publié en estimation communale tant
    qu'aucune source vectorielle n'est câblée. Le changer exige un amendement de
    la constitution du projet.
\end{itemize}

\begin{invariant}[title={Deux seuils, jamais confondus}]
Un seuil de 0,3 m gouverne les \emph{dommages} (compteurs d'objets touchés) ;
la carte, elle, représente l'eau dès la première profondeur positive. Confondre
les deux ferait disparaître de la carte des zones réellement inondées.
\end{invariant}

\section{Vérification de l'installation}

\begin{lstlisting}[language=bash,caption={Démarrage local et contrôle de vie.}]
# Backend
cd backend && ./venv/Scripts/python.exe -m uvicorn app.main:app --port 8000
# Frontend
cd frontend && npm run dev            # http://localhost:5173

# Controle de vie
curl http://127.0.0.1:8000/health     # {"status":"ok"}

# Suite de tests hors-ligne
cd backend && ./venv/Scripts/python.exe -m pytest tests/ -q

# Compilation de ce rapport
tectonic -X compile rapport/main.tex
\end{lstlisting}


\end{document}
